\section{Related Work}

\subsection{Adversarial benchmarks for T2I safety}
A growing body of work stress-tests safety mechanisms in generative models by crafting prompts that elicit unsafe outputs. Ring-A-Bell~\citep{tsai2023ring}, UnlearnDiff~\citep{zhang2024generate}, and MMA-Diffusion~\citep{yang2024mma} attack nudity generation with progressively more implicit phrasing, and P4D-N~\citep{chin2023prompting4debugging} augments these with automated prompt optimization against existing defenses. A recent line extends the attack surface beyond keyword-based nudity: Metaphor-based Jailbreaking Attack (MJA)~\citep{zhang2026reason2attack} uses LLM reasoning to produce prompts that evade keyword filters entirely (\eg, ``showing raw strength'' for violence) while still eliciting visually recognizable harmful patterns, and currently exposes the largest gap between training-free defenses and ground truth. Our evaluation unifies these four nudity benchmarks with I2P~\citep{schramowski2023safe} for multi-concept coverage and the four MJA categories for metaphorical adversarial pressure, and connects per-concept difficulty to detection modality through a probe-dominance analysis (Appendix~\ref{app:probe_dominance}).

\subsection{Safe T2I generation}

\textbf{Training-based.} A prominent line of work modifies diffusion model weights to remove unsafe concepts. ESD~\citep{gandikota2023erasing} fine-tunes cross-attention layers to minimize the probability of target concepts; SDD~\citep{kim2023towards} uses self-distillation to steer weights toward a safer distribution; UCE~\citep{gandikota2024unified} provides closed-form weight edits that remove the iterative training loop; and RECE~\citep{gong2024reliable} strengthens reliability through eraser ensembles. Concept Ablation~\citep{kumari2023ablating} and Forget-Me-Not~\citep{zhang2024forget} offer alternative fine-tuning objectives. While effective at per-concept suppression, these methods require separate training per target, risk catastrophic forgetting on nearby benign concepts, and remain constrained even on safe prompts; they also do not adapt cheaply when the adversarial distribution shifts (\eg, moving from keyword-based to metaphor-based attacks typically demands a fresh round of fine-tuning). In contrast, \ours{} is training-free, adapts to new concepts by dropping in a handful of exemplars, and leaves the base model fully intact on benign prompts.

\textbf{Training-free.} An alternative paradigm preserves model weights and intervenes only during inference. SLD~\citep{schramowski2023safe} manipulates classifier-free guidance at four preset strengths to steer denoising away from harmful concepts; SEGA~\citep{brack2023sega} extends CFG with semantic-guidance directions; SAFREE~\citep{yoon2024safree} projects harmful tokens orthogonal to a toxic subspace in joint text and latent spaces; and Safe Denoiser~\citep{kim2025training} derives a safer sampling trajectory from a negation-set formulation and can be stacked on top of SAFREE as a compound baseline. These methods share a common weakness: suppression is applied \emph{globally} across all spatial locations and timesteps, inducing semantic drift and quality degradation on the benign part of a prompt. \ours{} instead replaces a single global direction with a family-indexed exemplar pack, adds a pixel-wise winner-take-all partition that intervenes only on detected regions, and admits two complementary How modes (anchor inpainting and hybrid guidance) within a single framework.

\subsection{Safety evaluators for T2I}
Evaluating a safety mechanism honestly is as delicate as designing one. Automated classifiers such as NudeNet~\citep{notai_nudenet_2019} and Q16~\citep{schramowski2022can} answer a single harmful-or-not question and therefore implicitly count images driven into blur, noise, or semantic collapse as \emph{safe}, conflating collapse with erasure. VLM-based evaluators return finer-grained judgments, and MJ-Bench~\citep{chen2024mj} provides a standardized protocol for benchmarking such judges on safety-preference tasks. We propose a four-class content-based rubric, \{\textsc{Safe}, \textsc{Partial}, \textsc{Full}, \textsc{NotPeople}/\textsc{NotRelevant}\}, that explicitly flags semantic collapse rather than counting it as safe, and instantiate it with Qwen3-VL-8B~\citep{bai2025qwen3} as the safety evaluator. The rubric is validated on MJ-Bench and further confirmed against a human-majority annotation study (Appendix~\ref{app:human_study}); the full validator analysis is in Appendix~\ref{app:evaluator}.
